% DFSLICE1
\documentclass[11pt]{article}
\usepackage[margin=.88in]{geometry}
\usepackage{amsmath,amssymb}
\begin{document}
\begin{center}
{\Large DFSLICE1: exchangeability does not supply the $\varepsilon$-box}\\[4pt]
\textbf{Date: 2026-10-04 22:27 PT}\\
\textbf{Status: literature comparison + COMPUTATIONAL check}\\
\textbf{P versus NP status: OPEN}
\end{center}

Diaconis--Freedman (Ann.\ Probab.\ 8, 1980) bound the total variation between sampling with and without replacement by $2ck/n$ for an alphabet of size $c$, and show that the rate $k(k-1)/n$ is sharp for an infinite alphabet. On a Boolean slice, $c=2$, so the with-replacement bound is $4r/M$.

TR26-220, Definition 2.2, asks for something stricter. A marginal on at most $k$ coordinates must be a mixture of product laws whose biases all lie in $[p(1-\varepsilon),p(1+\varepsilon)]$. The method-of-moments argument in their Section 2 then centers $X_i=Y_i-\alpha$ and partitions $\mathbb E[X_{i_1}\cdots X_{i_p}]$ by the dependency graph of the terms. Exchangeability does not put those biases in the box.

The measured total variation of the exact density-$3.5$ slice against independent Bernoulli$(t/M)$ sits well below the Diaconis--Freedman number:
\[
\begin{array}{c|c|c|c}
q&r&\mathrm{TV}&4r/M\\
\hline
20&2&1.67\times 10^{-6}&8.77\times 10^{-4}\\
20&4&9.80\times 10^{-6}&1.75\times 10^{-3}\\
20&8&4.37\times 10^{-5}&3.51\times 10^{-3}
\end{array}
\]
The same pattern holds at $q=16$ and $q=24$. The exchangeability bound is valid and loose. It is not the obstruction.

The binding constraint remains SLICE-PSEUDO-LOWER: $\varepsilon^2\ge (k-1)(M-t)/(t(M-1))$. That is what closes the TR26-220 hypothesis for width at least $4$ under the budget $\varepsilon\le 1/w$. A small total-variation gap does not reopen it. Both facts are generic to every fixed-weight slice, including the GE-easy XOR slice. $P\ne NP$ remains OPEN.
\end{document}
